\begingroup

\setlength{\LTcapwidth}{\textwidth}

\begin{longtable}{@{}>{\raggedright\arraybackslash}p{1.400cm}>{\raggedleft\arraybackslash}p{2.200cm}>{\raggedright\arraybackslash}p{3.000cm}>{\raggedleft\arraybackslash}p{3.400cm}@{}}

\caption{Ringkasan hasil System Usability Scale pada dua belas responden.}

\label{tab:tab6.19_ringkasan_sus} \\

\toprule

\textbf{Butir} & \textbf{Rerata} & \textbf{Arah} & \textbf{Sumbangan SUS} \\

\midrule

\endfirsthead



\multicolumn{4}{@{}l}{\small\emph{Tabel \thetable{} (lanjutan)}} \\

\toprule

\textbf{Butir} & \textbf{Rerata} & \textbf{Arah} & \textbf{Sumbangan SUS} \\

\midrule

\endhead



\bottomrule

\endlastfoot



1 & 3,750 & Positif & 6,88 \\



2 & 3,250 & Negatif & 4,38 \\



3 & 3,750 & Positif & 6,88 \\



4 & 3,917 & Negatif & 2,71 \\



5 & 4,083 & Positif & 7,71 \\



6 & 3,000 & Negatif & 5,00 \\



7 & 3,750 & Positif & 6,88 \\



8 & 3,000 & Negatif & 5,00 \\



9 & 3,833 & Positif & 7,08 \\



10 & 3,917 & Negatif & 2,71 \\

\end{longtable}



\noindent\tabnotestyle Skala butir 1--5. Rerata SUS = 55,21 / 100 (SB 9,20; median 51,25; rentang 40,0--72,5; selang kepercayaan 95\% 49,36--61,05). Skor dihitung per responden dengan algoritma Brooke dari $n = 12$ baris tanggapan, sehingga sebaran antarresponden teridentifikasi.\normalsize

\endgroup

